\section{The uniform corona problem}
Let $M$ be a Riemann surface of genus $g\geq 2$. It is well known 
that the coordinate ring $\oplus_{m=0}^\infty H^0(M,K_M^m)$
is finitely generated. That is, there is an $m_0>0$ such that for any
$m>0$ and $S\in H^0(M,K_M^m)$, $S$ can be represented by
\begin{equation}\label{corona}
S=\sum_{i=1}^{d} U_iT_i,
\end{equation}
where $U_i\in H^0(M,K_M^{m_0})$ and $T_i\in H^0(M,K_M^{m-m_0})$
for $i=1,\cdots, d=\dim H^0(M,K_M^{m_0})$.
Finding a suitable set of $\{T_i\}_{i=1,\cdots,d}$ is called the 
corona problem(cf. ~\cite{JBG}).



We need to consider the case where
$M$ approaches to the boundary of the moduli space
in the Teichm\"uller theory. So
in addition to the existence of $U_i$ and $T_i$, we need some uniform 
estimates. 
In this section, we give the uniform estimate for the corona problem on
Riemann surfaces.


\begin{theorem}\label{5main}
Let $M$ be a Riemann surface 
of genus $g$ and constant curvature $(-1)$.
Then there is an $m_0>0$ such that for any $m>m_0$ and $S\in
H^0(M,K_M^m)$,
there is a decomposition
\[
S=\sum_{i=1}^d S_i
\]
of $S_i\in H^0(M,K_M^m) (i=1,\cdots, d)$
such that
\begin{align}\label{result}
\begin{split}
&||S_i||_{L^2}\leq C(m,m_0,g)||S||_{L^2}\\
&||S_i||_{L^\infty}\leq C(m,m_0,g)||S||_{L^\infty}
\end{split}
\end{align}
for $i=1,\cdots d$,
and 
\[
S_i=T_iU_i
\]
for a basis $U_1,\cdots,U_d$ of $H^0(M,K_M^{m_0})$
and $T_1,\cdots,T_d\in H^0(M,K_M^{m-m_0})$.
\end{theorem}

\begin{rem}
An estimate on $T_i (i=1,\cdots, d)$ alone is not expected because of the
counterexample in \S 3, where $\sum||U_i||^2$ can be arbitrarily small. 
\end{rem}


Throughout this section, we will use the notation $D_1, D_2,\cdots$ to
denote the
constants depending only on $m_0$ and the genus $g$.
We also use $A\leq\sim B$ to mean that there is a
positive constant $C=C(m_0,
g)$,
depending only on $m_0$ and $g$ such that $A\leq CB$. Likewise, we use
$A\geq\sim B$ to denote the fact $A\geq CB$
for some constant $C=C(m_0, g)$.


The idea of the proof is that, if the injective radius of $M$ is greater
than
an absolute  constant, then $\sum||U_i||^2$ has a lower bound by an
absolute
 positive constant. In this case, we can solve the corona problem
exactly using the method in~\cite{JBG}. So  we just need to prove the
theorem in the case where $inj(M)$ is arbitrarily small. By the collar
theorem, we know that in this case, there are finite many collars
$C_{\delta_1},\cdots, C_{\delta_s}$ 
(with $\max\delta_i$ small)
embedded in $M$ and they do not
intersect each other. By Proposition~\ref{prop31}, outside the collars,
the injective radius has an absolute lower bound. Special
care must be taken for the sections over these collars.
In order to take care of the collars to get the estimates, 
we first fix a
collar $C_\delta$ embedded in $M$ with the parameter $\delta$
small. We will use all the notations about collars in \S~\ref{sec3}.
For any $\tilde R>0$, 
let
\[
C_\delta(\tilde R)=\{x||\rho|\leq\tilde R\}.
\]
In particular, $C_\delta=C_\delta(R)$ with
$\delta\sinh R=\eps_1$.

We choose and fix a number $m_0>0$ such that $K_M^{m_0}$ is very ample. 
Let $\tilde\eta$ be the cut-off function of $M$ defined as
\begin{equation}\label{etai}
\tilde\eta=
\left\{
\begin{array}{ll}
\eta(\rho-(R-1)) &\rho\geq 0\quad and\quad x\in C_\delta,\\
\eta(-\rho-(R-1)) & \rho<0\quad and\quad x\in C_\delta,\\
0 & otherwise,
\end{array}
\right.
\end{equation}
where the function $\eta$ is defined in ~\eqref{eta}.
Let
\begin{equation}\label{5-u1}
u_1=\frac{1}{\sqrt\delta}\tilde \eta(dz)^{m_0}
\end{equation}
be a section of $K_M^{m_0}$ over $M$ using this cut-off function,
where $dz$ is defined
in~\eqref{dz}.
We can check that
\[
\left\{
\begin{array}{ll}
||\bar\pa u_1||^2\leq\frac{4}{\delta(\cosh\rho)^{2m_0}} &
R-1\leq|\rho|\leq
R,\\
\bar\pa u_1=0 & otherwise.
\end{array}
\right.
\]
 Thus
\[
\int_M||\bar\pa u_1||^2\leq 2\int_{R-1}^R
\frac{4}{\delta(\cosh\rho)^{2m_0}}\cosh\rho\delta d\rho
\leq \frac{8}{(\cosh (R-1))^{2{m_0}-1}}.
\]
By Proposition~\ref{del}, there is a section $u$ of $K_M^{m_0}$ such that
\[
\bar\pa u=\bar\pa u_1
\]
with 
\[
\int_M||u||^2\leq\frac{1}{m_0-1}\int_M||\bar\pa u_1||^2
\leq\frac{8}{(m_0-1)(\cosh (R-1))^{2m_0-1}}.
\]
By using ~\eqref{eps2}, we see that
\begin{equation}\label{5-u3}
\int_M ||u||^2\leq\sim\delta^{2m_0-1}.
\end{equation}
Let 
\begin{equation}\label{5-defu}
U'=u_1-u.
\end{equation}
Then $\bar\pa U'=0$. 


\begin{lemma}\label{claim}
Let $U''$ be a holomorphic section of $K_M^{m_0}$ 
on $C_\delta(R)$
such that
$(U'',u_1)_{C_\delta(R-2)}=0$. That is,
\begin{equation}\label{prep}
\int_{C_\delta(R-2)} <U'',u_1>=0,
\end{equation}
where $u_1$ is the section defined in~\eqref{5-u1}. Then there
is an absolute constant $\eps_6>0$ such that
\begin{equation}\label{5-n1}
||U''||\leq\sim e^{-\eps_6e^{R-|\rho|}+m_0(R-|\rho|)}||U''||_{L^2
(C_\delta)},
\end{equation}
and
\begin{equation}\label{5-n2}
||\nabla U''||\leq
\sim e^{-\eps_6e^{R-|\rho|}+(m_0+1)(R-|\rho|)}||U''||_{L^2
(C_\delta)}.
\end{equation}
\end{lemma}

We use the notation in \S ~\ref{sec3}. Let
\begin{align*}
& w_{1}=e^{\frac{2\pi\sqrt{-1}}{\delta}\theta
-\frac{4\pi}{\delta}
(\arctan e^{\rho}-\arctan e^{-(R-2)})},\\
& w_{2}=e^{-\frac{2\pi\sqrt{-1}}{\delta}\theta
+\frac{4\pi}{\delta}
(\arctan e^{\rho}-\arctan e^{R-2})}.
\end{align*}
Let   
\[
U''=(g_1(w_1)+a+g_3(w_2))(dz)^{m_0}
\]
be the decomposition similar to that in~\eqref{g} where
$g_1(w_1)$ and $g_2(w_2)$ are holomorphic functions of $w_1$ and
$w_2$ respectively, $g_1(0)=g_2(0)=0$, $a$ is a constant,
and $dz$ is defined in~\eqref{dz}.

Using~\eqref{prep}, we see that $a=0$. By the Schwartz Lemma, we have
\begin{equation}\label{sch}
|g_1(w_1)|\leq |w_1| \max_{|w_1|=1}|g_1(w_1)|.
\end{equation}
At each point of $\{|w_1|=1\}$ or $\{\rho=R-2\}$, by the collar theorem,
there is an
absolute lower bound for the injective radius. Thus by the 
Cauchy integral formula,
we have
\begin{equation}\label{5-8}
\max_{|w_1|=1}||g_1(dz)^{m_0}||\leq
\sim||g_1(dz)^{m_0}||_{L^2(C_\delta)}.
\end{equation}

Using ~\eqref{sch} and~\eqref{5-8}, we have
\begin{equation}\label{5-8-1}
||g_1 (dz)^{m_0}||\leq e^{-\frac{4\pi}{\delta}(\arctan e^\rho-\arctan
e^{-(R-2)})}
\frac{(\cosh R)^{m_0}}{(\cosh\rho)^{m_0}}
||g_1(dz)^{m_0}||_{L^2(C_\delta)}.
\end{equation}
It is elementary to check that there is an absolute constant
$\eps_6>0$ such that
\begin{equation}\label{5-8-1-1}
e^{-\frac{4\pi}{\delta}(\arctan e^\rho-\arctan e^{-(R-2)})}
\leq e^{-\eps_6e^{R-|\rho|}}.
\end{equation}
Combining~\eqref{5-8-1} and~\eqref{5-8-1-1}, 
\begin{equation}\label{5-13a}
||g_1(dz)^{m_0}||\leq
\sim e^{-\eps_6e^{R-|\rho|}+m_0(R-|\rho|)}
||g_1(dz)^{m_0}||_{L^2(C_\delta)}.
\end{equation}
Similarly, we have
\begin{equation}\label{5-14a}
||g_2(dz)^{m_0}||\leq
\sim e^{-\eps_6e^{R-|\rho|}+m_0(R-|\rho|)}
||g_2(dz)^{m_0}||_{L^2(C_\delta)}.
\end{equation}
~\eqref{5-13a} and
~\eqref{5-14a} give the inequality~\eqref{5-n1}.

On the other hand, a straightforward
computations gives
\begin{align}
\begin{split}
&\nabla U''=(\frac{2\pi\sqrt{-1}}{\delta}(w_1 g_1'+w_2
g_2')\\
&\qquad -\frac 12 m_0\sqrt{-1}\sinh\rho (g_1+ g_2))
(dz)^{m_0}\otimes dz.
\end{split}
\end{align}
Using the same argument as above, we get~\eqref{5-n2}.

\qed



\begin{lemma}\label{lem51}
With the notations as above, 
there is a  constant $r>2$, depending only on $m_0$ and the genus
$g$, such
that
 $U'\neq 0$ on $C_\delta(R-r)$, where $U'$ is defined in~\eqref{5-defu}.
\end{lemma}


{\bf Proof.}
By~\eqref{5-defu}, we know that $u$ is holomorphic on $C_\delta(R-2)$. Let
\begin{equation}\label{alpha1}
u=u'+\alpha u_1
\end{equation}
be the decomposition of $u$ such that $(u',u_1)
_{C_\delta(R-2)}=0$ and $\alpha$ is
a constant.
Then
\begin{equation}\label{5-5}
\int_M||u||^2\geq\int_{C_{\delta}(R-2)}
||u'||^2
+\int_{C_{\delta}(R-2)}
|\alpha|^2||u_1||^2.
\end{equation}
In particular
\begin{equation}\label{5-6}
 \int_{C_{\delta}(R-2)}
||u||^2
\geq |\alpha|^2\int^{R-2}_{-(R-2)}\frac{1}{(\cosh\rho)^{2m_0-1}}
d\rho.
\end{equation}
By~\eqref{5-u3} and ~\eqref{5-6}, we have
\begin{equation}\label{5-7}
|\alpha|
\leq\sim \delta^{m_0-\frac 12}.
\end{equation}
On the other hand, by Lemma~\ref{claim},
\begin{equation}\label{5-7-3}
||u'||\leq\sim
e^{-\eps_6e^{R-|\rho|}+m_0(R-|\rho|)}\delta^{m_0-\frac 12}.
\end{equation}
Thus by~\eqref{5-u1}, ~\eqref{5-7} and~\eqref{5-7-3}, there are constants
$D_1$ and $D_2$ such that
\begin{align}\label{5-7-4}
\begin{split}
&||U'(x)||\geq\frac{1}{\sqrt\delta(\cosh\rho)^{m_0}}
(1-D_1\delta^{m_0-\frac 12})\\
&-D_2
e^{-\eps_6e^{R-|\rho|}+m_0(R-|\rho|)}\delta^{m_0-\frac 12}
\end{split}
\end{align}
for $|\rho|<R-3$.
If $r$ is large enough, then $||U'(x)||>0$ for $|\rho|<R-r$.
In particular $U'\neq 0$ on $C_\delta(R-r)$. This completes the proof of
the lemma.

\qed


Now we assume the general case. Let $C_{\delta_1},\cdots, C_{\delta_s}$ be
the collars of parameters $\delta_1,\cdots, \delta_s$ respectively
embedded
into $M$. We assume that $C_{\delta_i}, (i=1,\cdots, s)$ do not intersect
each other.
By Lemma~\ref{lem51}, there are $R_i, (i=1,\cdots, s)$ such that we can
find
$U_i'\in H^0(M,K_m^{m_0}), (i=1,\cdots, s)$ with
$U_i'|_{C_{\delta_i}(R_i)}\neq 0$. Let $\rho_i, dz_i, (i=1,\cdots, s)$ be
defined in \S~\ref{sec3}
coresponding to $C_{\delta_i} (i=1,\cdots, s)$, respectively. Let
\[
\eta_i=
\left\{
\begin{array}{ll}
\eta(\rho_i-(R_i-1)) &\rho_i\geq 0\quad and\quad x\in C_{\delta_i}\\
\eta(-\rho_i-(R_i-1)) &\rho_i\leq 0\quad and\quad x\in C_{\delta_i}\\
0 & otherwise
\end{array}
\right.
\]
be the cut-off functions for $i=1,\cdots, s$. Assume that $\max \delta_i$
is small enough.
Then
by Lemma~\ref{lem51}, $U_1',\cdots, U_s'\in H^0(M,K_M^{m_0})$
have  the following properties 
\begin{enumerate}
\item $||U_i'||\geq\sim\frac{1}{\sqrt{\delta_i}}e^{-m_0|\rho_i|}$ on
$C_{\delta_i}(R_i)$ for $i=1,\cdots, s$ 
(by ~\eqref{5-7-4});
\item
There are decompositions 
 $U_i'|_{C_{\delta_i}(R_i)}=\alpha_iv_{1i}+v_{2i}$ with
\begin{equation}\label{vj}
v_{1i}=\frac{1}{\sqrt{\delta_i}}\eta_i(dz_i)^{m_0},\qquad
1\leq i\leq s
\end{equation}
and
\[
(v_{1i},v_{2i})_{C_{\delta_i}(R_i)}=0, \qquad 1\leq i\leq s,
\]
where $\alpha_i (i=1,\cdots,s)$ are a constant such that for
$1\leq i\leq s$,
\begin{equation}\label{vj1}
1\leq\sim\alpha_i\leq\sim 1;
\end{equation}
\item By~\eqref{5-u1} and~\eqref{5-u3},
\begin{equation}\label{5-7-5}
\int_{M\backslash C_{\delta_i(R_i)}}||U_i'||^2\leq\sim \delta_i^{2m_0-1}
\end{equation}
for $i=1,\cdots, s$.
\end{enumerate}

We have the following lemma:

\begin{lemma}\label{lem531}
With the notations as above,
there are holomorphic sections $U_1,\cdots, U_s\in 
H^0(M,K_M^{m_0})$ such that
\begin{enumerate}
\item $||U_i||\geq\sim \frac{1}{\sqrt\delta_i}e^{-m_0|\rho_i|}$ on 
$C_{\delta_i}(R_i)$;
\item $(U_i|_{C_{\delta_j}(R_j)},v_{1j})_{C_{\delta_j}(R_j)}=0$ for $i\neq
j, 1\leq i,j\leq
s$;
\item $\int_{M\backslash
C_{\delta_i}(R_i)}||U_i||^2\leq\sim
\delta_i^{2m_0-1}$.
\end{enumerate}
\end{lemma}

{\bf Proof.} Let
\[
\beta_{ij}=(U_i',v_{1j})_{C_{\delta_j}(R_j)}, \qquad 1\leq i,j\leq s.
\]
Then if $i\neq j$ we have
\begin{equation}\label{b-1}
|\beta_{ij}|\leq\sim\delta_i^{m_0-\frac 12}
\end{equation}
by~\eqref{5-7-5} and the definition of $v_{1j},
(j=1,\cdots, s)$ in~\eqref{vj}. We also have
\begin{equation}\label{b-2}
1\geq\sim 
\beta_{ii}\geq\sim 1
\end{equation}
by~\eqref{vj1}. 

Let $B=(\beta_{ij})_{s\times s}$ be the matrix of $(\beta_{ij})$ for
$1\leq i,j\leq s$. Then   by~\eqref{b-1}, ~\eqref{b-2}, $B$  is an
invertible matrix,
when $\max\delta_i$ is small enough.
Let $A=B^{-1}$ be the
inverse matrix and let $A=(\alpha_{ij})_{s\times s}$. Define
\[
U_i=\sum_{j=1}^s\alpha_{ij}U_j', \quad 1\leq i\leq s.
\]
Then $U_i (i=1,\cdots, s)$ satisfies all the requirements in the lemma
by the fact that
\[
\begin{array}{ll}
|\alpha_{ij}|\leq\sim \delta_i^{m_0-\frac 12} & i\neq j,\\
1\geq\sim \alpha_{ij}\geq\sim 1 & i=j.
\end{array}
\]


\qed

Let $U_{s+1}',\cdots, U_d'$ be an orthonormal basis of the space 
$span\{ U_1,\cdots. U_s\}^\perp$. Assume that
\[
(U_i', U_j)=0
\]
for $s<i\leq d, 1\leq j\leq s$.
Let
\begin{equation}\label{kk}
U_i=U_i'-\sum_{j=1}^s\gamma_{ij}U_j,\quad s<i\leq d,
\end{equation}
where
\begin{equation}\label{b-3}
\gamma_{ik}=\frac{1}{(U_k, v_{1k})_{C_{\delta_k}(R_k)}}
(U_i', v_{1k})_{C_{\delta_k}(R_k)},
\quad 1\leq k\leq s, s<i\leq d.
\end{equation}
Then we have
\[
(U_i|_{C_{\delta_j}(R_j)}, v_{1j})_{C_{\delta_j}(R_j)}=0
\]
for $s<i\leq d$ and $1\leq j\leq s$.


We have the following lemma:









\begin{lemma}\label{lem53}
If $x\notin \cup_{j=1}^sC_{\delta_j}(R_j-2)$, then 
\[
\sum_{i=1}^d||U_i||^2\geq\sim 1.
\]
\end{lemma}


{\bf Proof.} Let
\[
l_{ij}=(U_i, U_j)
\]
for $1\leq i,j\leq d$. If $1\leq i,j\leq s$, we have
\begin{equation}\label{k1}
\begin{array}{ll}
(U_i, U_j)\leq\sim \max \delta_i & i\neq j,\\
1\geq\sim (U_i, U_j)\geq\sim  1 & i=j
\end{array}
\end{equation}
by~\eqref{b-1}, ~\eqref{b-2} and lemma~\ref{lem531}. 

If $1\leq i\leq s, s<j\leq d$, we have
\[
(U_i, U_j)=-\sum_{k=1}^s\gamma_{jk} (U_i, U_k)
\]
using ~\eqref{kk}.
By the definition of $\gamma_{jk}$ in~\eqref{b-3}, we have
\begin{equation}\label{g1}
|\gamma_{jk}|\leq\sim\max\delta_j.
\end{equation}
Thus
\begin{equation}\label{k2}
|(U_i, U_j)|\leq\sim\max\delta_j
\end{equation}
for $1\leq i\leq s$ and $s<j\leq d$.
Finally, if $s<i,j\leq d$, then 
\begin{equation}\label{k3}
\begin{array}{ll}
|(U_i, U_j)|\leq\sim \max\delta_i & i\neq j,\\
1\geq\sim (U_i, U_j)\geq\sim 1 & i=j
\end{array}
\end{equation}
by~\eqref{kk} and~\eqref{g1}.
Using~\eqref{k1}, ~\eqref{k2} and~\eqref{k3}, we have
\[
\begin{array}{ll}
|l_{ij}|\leq\sim \max \delta_i & i\neq j,\\
1\geq\sim |l_{ij}|\geq\sim 1 & i=j.
\end{array}
\]

Let $(m_{ij})_{d\times d}$ be the matrix such that
\[
\sum_{i=1}^d\sum_{t=1}^d m_{ji}\bar{m_{tk}}l_{it}=\delta_{jk}
\]
for $1\leq j,k\leq d$. We can choose $m_{ij}$ such that
\begin{equation}\label{m-1}
\begin{array}{ll}
|m_{ij}|\leq\sim \max
\delta_i & i\neq j,\\
1\geq\sim m_{ij}\geq\sim 1 & i=j.
\end{array}
\end{equation}

It is easy to check that
$\sum_{j=1}^d m_{ij} U_j$ for $1\leq i\leq d$ forms an orthonormal basis
of $H^0(M, K_M^{m_0})$.
Thus we have
\[
\sum_{i=1}^d||\sum_{j=1}^dm_{ij}U_j||^2\geq\sim 1
\]
by Theorem~\ref{third}.
The lemma thus follows from~\eqref{m-1} and the fact that $\max \delta_i$
is
small.

\qed



We summarize the results up to now in the following
\begin{prop}\label{prop51}
Let $U_1,\cdots, U_d$ be sections of $H^0(M,K_M^{m_0})$
as above. Then
\begin{align}
&||U_i||+||\nabla U_i||\leq\sim 1 &x\notin
\cup_{j=1}^sC_{\delta_j}(R_j-2),
\label{5-88a}\\
&||U_i||+||\nabla U_i||\leq\sim
e^{-(m_0+1)(R_j-|\rho_j|)}
&x\in C_{\delta_j}(R_j-2), 1\leq j\leq s, \label{5-88b}
\end{align}
for $i>s$ and
\begin{align}
&||U_i||+||\nabla U_i||\leq\sim
{\delta_i}^{m_0-\frac 12}
&x\notin\cup_{j=1}^s C_{\delta_j}(R_j-2),\label{5-88c}\\
&||U_i||+||\nabla U_i||\leq\sim
\frac{1}{\sqrt\delta_i e^{m_0|\rho_i|}}
&x\in C_{\delta_i}(R_i-2),\label{5-88d}\\
&||U_i||+||\nabla U_i||\leq\sim\delta_i^{m_0-\frac 12}
e^{-(m_0+1)(R_j-|\rho_j|)} &
x\in C_{\delta_j}(R_j-2), j\neq i.\label{5-88e}
\end{align}
for $1\leq i\leq s$.
Furthermore, 
\begin{equation}\label{new}
\sum_{k=1}^d||U_i||^2\geq\sim 1
\end{equation}
for $x\notin \cup_{i=1}^s C_{\delta_i}(R_i-2)$.
\end{prop}

{\bf Proof.}  If $x\notin \cup_{j=1}^sC_{\delta_j}(R_j)$, then $\delta_x$
has a uniform lower bound.
~\eqref{5-88a} follows from the Cauchy integral formula. ~\eqref{5-88b}
is a corollary of Lemma~\ref{claim}. ~\eqref{5-88c} follows from
~\eqref{5-u3} and the Cauchy integral formula. ~\eqref{5-88d} follows
from
a straightforward computation.~\eqref{5-88e} follows from
Lemma~\ref{claim} and~\eqref{5-7-5}. Finally, ~\eqref{new} is just a
restatement of the conclusion of Lemma~\ref{lem53}.

\qed



Define an inner product $<\,,\,>$ in the coordinate ring $\oplus_{m=0}^
\infty H^0(M,K_M^m)$. Let $S_1\in H^0(M,K_M^m)$ and $S_2\in
H^0(M,K_M^{m_1})$. Suppose that $m\ge m_1$. We define a section of
$K_M^{m-m_1}$ as follows: Let $x\in M$ and $U_x$ is a local trivialization
of 
$K_M$. let $S_1|_{U_x}=S'\cdot S''$ for 
$S'\in \Gamma(U_x, K_M^{m-m_1})$ and $S''\in
\Gamma(U_x, K_M^{m_1})$. Then
\[
S_3|_{U_x}=S'<S'', S_2>_{H^{m_1}},
\]
where $<,\,,\,>_{H^{m_1}}$ is the pointwise inner product.


{\bf Proof of Theorem~\ref{5main}.} 
We modify the method of Wolff's~\cite{WT} of solving the corona problem
on the unit disk.
 First we construct a $C^\infty$ solution. Let $S\in
H^0(M,K_M^m)$ for a fixed $m>m_0$.
Let
\begin{equation}\label{5-b1}
\left\{
\begin{array}{ll}
 b_k=\eta_k\frac{S}{U_k}+(1-\sum_{j=1}^s\eta_j)
\frac{<S, U_k>}{\sum_{j=1}^d||U_j||^2}& 1\leq k\leq s,\\
 b_k=(1-\sum_{j=1}^s\eta_j)\frac{<S, U_k>}{\sum_{j=1}^d||U_j||^2}
& k>s.
\end{array}
\right.
\end{equation}
Here $b_k (1\leq k\leq s)$ is well defined
because of Lemma~\ref{lem51}. We can check that
\[
S=\sum_{k=1}^d U_k b_k.
\]



If $x\notin \cup_{j=1}^sC_{\delta_j}(R_j-2)$, then by Lemma~\ref{lem53},
\begin{equation}\label{5-b2}
\left\{
\begin{array}{ll}
||b_k||\leq\sim
(\sqrt\delta_k (\cosh\rho_k)^{m_0}+1)||S||& 1\leq k\leq s,\\
 ||b_k||\leq\sim ||S||
&  k>s.
\end{array}
\right.
\end{equation}
By~\eqref{5-b1}, we have
\begin{align}\label{5-b3}
\begin{split}
&\bar\pa b_k=\bar\pa\eta_k\frac{S}{U_k}
-\bar\pa\sum_{j=1}^s\eta_j
\frac{<S,U_k>}{\sum_{j=1}^d||U_j||^2}
+(1-\sum_{j=1}^s\eta_j)\frac{<S,\nabla U_k>}
{\sum_{j=1}^d||U_j||^2}\\\quad-
&(1-\sum_{j=1}^s\eta_j)\frac{<S,U_k>\sum_{j=1}^d<U_j,\nabla U_j>}
{(\sum_{j=1}^d||U_j||^2)^2}\qquad 1\leq k\leq s,\\
&\bar\pa b_k=
-\bar\pa\sum_{j=1}^s\eta_j
\frac{<S,U_k>}{\sum_{j=1}^d||U_j||^2}
+(1-\sum_{j=1}^s\eta_j)\frac{<S,\nabla U_k>}
{\sum_{j=1}^d||U_j||^2}\\
&\quad-
(1-\sum_{j=1}^s\eta_j)\frac{<S,U_k>\sum_{j=1}^d<U_j,\nabla U_j>}
{(\sum_{j=1}^d||U_j||^2)^2}\quad  k>s.
\end{split}
\end{align}
Thus
\begin{equation}\label{5-b4}
\left\{
\begin{array}{ll}
||\bar\pa b_k||\leq \sim
\frac{1}{\delta_k^{m_0-\frac 12}}
||S||& 1\leq k\leq s,\\
||\bar\pa b_k||\leq\sim||S|| &  k>s.
\end{array}
\right.
\end{equation}
Let
\begin{equation}
c_{ik}=\frac{<\bar\pa b_k, U_i>}{\sum_{j=1}^d||U_j||^2}.
\end{equation}
for $i\neq k$ and
$1\leq i,k\leq d$. Then 
by~\eqref{5-b2} and~\eqref{5-b4}, we have
\begin{equation}\label{5-c1}
\left\{
\begin{array}{ll}
 ||c_{ik}||\leq\sim \frac{1}{\delta_k^{m_0-\frac 12}}
||S|| & 1\leq k\leq s,\\
 ||c_{ik}||\leq \sim 
||S|| &  k>s.
\end{array}
\right.
\end{equation}

Let's consider the equations
\[
\bar\pa b_{ik}=c_{ik}
\]
for $i\neq k$ and $1\leq i,k\leq d$. By Proposition~\ref{del}, the
solutions exist and we may assume that
\[
||b_{ik}||_{L^2}\leq ||c_{ik}||_{L^2}
\]
for $i\neq k, 1\leq i,k\leq d$. Thus
by~\eqref{5-c1}, we have
\begin{equation}\label{5-b5}
\left\{
\begin{array}{ll}
||b_{ik}||_{L^2}\leq\sim \frac{1}{\delta_k^{m_0-\frac 12}}
||S||_{L^2} & 1\leq k\leq s,\\
||b_{ik}||_{L^2}\leq \sim ||S||_{L^2}& 
 s<k\leq d.
\end{array}
\right.
\end{equation}

Let
\[
T_i=b_i+\sum_{k=1}^d(b_{ik}-b_{ki}) U_k.
\]
One can check that $\bar\pa T_i=0$ and
\[
S=\sum_{i=1}^dT_iU_i.
\]

In order to prove the theorem, we need to estimate $||T_iU_i||$
for $1\leq i\leq d$. By~\eqref{5-b5} and Proposition~\ref{prop51}, we have
\begin{equation}\label{5-e-1}
||T_iU_i||_{L^2(M\backslash \cup_{j=1}^sC_\delta(R_j))}
\leq\sim ||S||_{L^2}.
\end{equation}
 By the Cauchy integral formula and Proposition~\ref{prop31}
\begin{equation}\label{5-e-2}
||T_iU_i||(x)\leq\sim ||S||_{L^2}
\end{equation}
for $1\leq i\leq d$ and $x\notin \cup_{j=1}^sC_{\delta_j}(R_j)$.

Let
\[
E_i(\rho)=\{x\in C_{\delta_i}(R_i)||\rho_i(x)|\geq\rho\}, \quad
1\leq i\leq s
\]
for positive number $\rho>0$.
By~\eqref{5-b5} and Proposition~\ref{prop51} again, we have
\begin{equation}\label{5-b6}
\left\{
\begin{array}{ll}
 ||T_iU_i-\eta_iS||_{L^2(E_j(\rho)-E_j(\rho+2))}\leq
\sim e^{-(R_i-\rho)}||S||_{L^2}& 1\leq i\leq s,\\
 ||T_iU_i||_{L^2(E_j(\rho)-E_j(\rho+2))}\leq\sim e^{-(R_i-\rho)}
||S||_{L^2} &i>s.
\end{array}
\right.
\end{equation}
 By the Cauchy formula,
\begin{equation}\label{5-b7}
\left\{
\begin{array}{l}
||T_iU_i-\eta_iS||(x)\leq\sim \frac{1}{\delta_x}e^{-(R_i-\rho)}
||S||_{L^2}\quad 1\leq i\leq s,\\
||T_iU_i||(x)\leq\sim \frac{1}{\delta_x}e^{-(R_i-\rho)}
||S||_{L^2}\quad i>s.
\end{array}
\right.
\end{equation}
for any $x\in E_j(\rho+\frac 12)-E_j(\rho+1)$.
Since
\[
\delta_x\geq\eps_5\delta_j(\cosh\rho_j)
\]
by~\eqref{ddelta}, we have
\begin{equation}\label{5-b8}
\left\{
\begin{array}{ll}
||T_iU_i-\eta_iS||(x)\leq\sim ||S||_{L^2} & 1\leq i\leq s,\\
||T_iU_i||(x)\leq\sim ||S||_{L^2} & i>s
\end{array}
\right.
\end{equation}
for any $x\in E_j(\rho+\frac 12)-E_j(\rho+1)$ and any $|\rho_j|<R_j-3$,
$j=1,\cdots s$. Thus
if $||S||_{L^2}=1$, then by~\eqref{5-e-2} and~\eqref{5-b8}
\[
||T_iU_i||_{L^2}\leq\sim 1
\]
for $1\leq i\leq d$. If $||S||_{L^\infty}=1$, then
\[
||T_iU_i||_{L^\infty}\leq\sim 1
\]
for $1\leq i\leq d$.
These results give the inequality  ~\eqref{result}.

\qed

